\documentclass[10pt,letterpaper]{article}
\usepackage{cornell-notes}

\title{Chapter 24: Triangular Inversion}
\author{}
\date{}

\setDocTitle{Chapter 24: Triangular Inversion}
\setDocSubtitle{Combinatorial Algorithms}
\setDocVersion{v1.0}
\setDocStatus{Study Companion}
\setDocOwner{Combinatorial Algorithms Cornell Notes}
\setDocAuthor{}
\setDocDate{\today}
\setDocSummary{Cornell-format study and review notes for Chapter 24: Triangular Inversion.}

\setCornellCollection{Combinatorial Algorithms Cornell Notes}
\setCornellUnitType{Chapter}
\setCornellUnitNumber{24}
\setCornellUnitTitle{Triangular Inversion}
\setCornellCompanionLabel{Study and review companion}

\hypersetup{pdftitle={Chapter 24: Triangular Inversion},pdfsubject={Cornell notes},pdfauthor={}}

\begin{document}
\maketitle
\begin{tcolorbox}[colback=Paper,colframe=Ink]{\small\bfseries\color{Teal}CORNELL NOTES}\par{\LARGE\bfseries Chapter 24: Invert a Triangular Array (INVERT)}\par\medskip\textbf{Topic:} Unit upper-triangular matrix inversion\hfill\textbf{Date:} \today\par\textbf{Study purpose:} Derive the recurrence and safe in-place order.\end{tcolorbox}
\begin{tcolorbox}[colback=Teal!5,colframe=Teal,title=Essential question]Which dependency order makes the inverse of a unit upper-triangular matrix computable by a simple recurrence, even in the same storage?\end{tcolorbox}
\begin{tcolorbox}[colback=white,colframe=Mist,title=Learning objectives]\begin{itemize}\item State the unit upper-triangular precondition.\item Derive inverse entries from $AB=I$.\item Identify the row and column dependency order.\item Explain why reverse column processing permits aliasing.\item State operation and numeric assumptions.\end{itemize}\textbf{Prerequisites:} matrix multiplication, triangular matrices, recurrences.\end{tcolorbox}
\section*{Cornell notes}
\begin{longtable}{@{}Q!{\color{Mist}\vrule width 1pt}A@{}}\cornellhead
\cue{What is the input class?}{$A$ is $n\times n$, has $a_{ii}=1$, and has $a_{ij}=0$ for $i>j$. Its inverse $B=A^{-1}$ is also unit upper triangular. No division is required because every diagonal pivot equals one.}
\cue{What recurrence follows from $AB=I$?}{For $i<j$,
\[
0=(AB)_{ij}=b_{ij}+\sum_{k=i+1}^{j}a_{ik}b_{kj},
\qquad
b_{ij}=-\sum_{k=i+1}^{j}a_{ik}b_{kj}.
\]
Also $b_{ii}=1$ and $b_{ij}=0$ below the diagonal.}
\cue{What is the calculation order?}{For a fixed column $j$, calculate rows $i=j-1,j-2,\ldots,1$. Every required $b_{kj}$ has $k>i$ and is already known. Process columns from $j=n$ down to $1$.}
\cue{Why can input and output alias?}{When columns are processed in reverse, entries of $A$ needed for future calculations have not yet been overwritten by their inverse values. The recurrence reads the current row to the right and inverse values already formed lower in the same column.}
\cue{Correctness invariant}{Before calculating $b_{ij}$, the suffix $b_{i+1,j},\ldots,b_{jj}$ satisfies the corresponding equations of $AB=I$. The recurrence chooses the unique value making equation $(i,j)$ zero. Induction across rows and columns establishes $AB=I$.}
\cue{Cost}{The number of multiply-add terms is cubic in $n$ with leading triangular-sum behavior; storage is the output matrix, and it may share storage with the input. The specialized routine is simpler than general inversion because no pivoting or division is needed.}
\cue{Cautions}{The diagonal must be exactly one for the stated recurrence. A general nonsingular triangular matrix requires division by $a_{ii}$. Integer arithmetic is exact only when the unit-triangular structure and values keep results representable.}
\cue{Retrieval practice}{\begin{enumerate}\item Why is $B$ upper triangular?\item What is $b_{ij}$ for $i<j$?\item Why move upward within a column?\item Why process columns backward?\item What changes for nonunit diagonal entries?\end{enumerate}}
\end{longtable}
\begin{tcolorbox}[colback=Ink!4,colframe=Ink,title=Cornell summary]
INVERT specializes matrix inversion to unit upper-triangular arrays. Expanding the equation $AB=I$ at position $(i,j)$ expresses $b_{ij}$ as the negative sum of already known entries lower in column $j$, while diagonal entries remain one and entries below the diagonal remain zero. Calculating upward within each column satisfies these dependencies. Processing columns from right to left also permits the inverse to overwrite the input matrix safely. The method needs no pivoting or division, but its contract does not extend unchanged to triangular matrices with arbitrary diagonal values.
\end{tcolorbox}
\end{document}
